\documentclass[10pt,a4paper,openany,twoside]{book}

\usepackage{layout_cours}


\begin{document}
\section{Bordas 2023}
\subsection{Exo 79 p.}
On veut montrer que $\vec{AE}=\vec{FC}$

Pour cela, on va exprimer ces deux vecteurs en fonction du vecteur $\vec{CD}$.

\bigskip
D'une part,
\begin{align*}
    \vec{AE}
    &= \frac{1}{4}\vec{AB}\\
    &= \frac{1}{4}\vec{DC} \quad \text{car ABCD est un parallélogramme}
\end{align*}

\bigskip
D'autre part, on sait que $\vec{DF} = \frac{3}{4}\vec{DC}$.

On veut faire apparaître le vecteur $\vec{FC}$. Pour cela on introduit la lettre $C$. D'après la relation de Chasles, $\vec{DF} = \vec{DC}+\vec{CF}$.

On a donc l'égalité $\vec{DC}+\vec{CF}= \frac{3}{4}\vec{DC}$

\medskip
En soustrayant $\vec{DC}$ de chaque côté de l'équation, on obtient:
$\vec{CF}= -\frac{1}{4}\vec{DC}$

Donc $\vec{FC} = \frac{1}{4}\vec{DC}$.

\bigskip
Finalement, on a $\vec{AE}=\vec{FC}$ et AECF est un parallélogramme.
\end{document}